\section{Results}

\subsection{OSDR contains an interlocking irradiation series, not a set of disjoint studies}

Of the \NumArabidopsisStudiesScreened{} \textit{Arabidopsis} studies in OSDR, most carry
no parsable time axis. Those that do fall into two families. The
$\gamma$-irradiation family yields \NumCohortATimedSamples{} timed samples across
\NumCohortATimepoints{} timepoints from 10~min to 72~h, plus
\NumCohortAAnchorSamples{} untreated samples that anchor $t=0$
(Table~\ref{tab:cohorts}).

The series interlocks rather than butt-joining. OSD-508 and OSD-510 share six timepoints
and OSD-498 shares four with OSD-508, giving \NumSeriesSharedTimepoints{} timepoints
measured by more than one accession --- enough to fit and remove a residual per-study
offset without extrapolation. This is a direct consequence of provenance: the three core
accessions are one experiment split across three deposits, not three independent
laboratories. We state this explicitly because it bounds the claim. The cross-study
element of cohort~A is OSD-782, and OSD-782 alone.

The spaceflight cohort is thinner: \NumCohortBSamples{} samples over only
\NumCohortBTimepoints{} timepoints (4, 6 and 8~d) across four ecotypes. It is reported as
a generalisation test, not as a co-equal result.

\begin{figure}[t]
\centering
\includegraphics[width=\linewidth]{../figures/fig1_cohort_design}
\caption{Pseudo-time-series coverage. Each marker is a timepoint measured by one
accession; marker area is the replicate count. Cohort A's accessions overlap at
\NumSeriesSharedTimepoints{} timepoints, which is what allows a residual between-study
offset to be fitted without extrapolation.}
\label{fig:design}
\end{figure}

\subsection{Batch harmonisation preserves the DNA-damage response and its genetic dependency}

After within-study referencing and offset removal, all \NumSentinelsResponding{} canonical
DNA-damage sentinel genes exceed 2-fold induction in the wild-type arm, tracing a clean
rise--peak--decay kinetic: SMR7, for example, rises from $1.34$ at 10~min to a maximum of
$7.62$ ($\log_2$) at 3~h before decaying to $4.42$ at 24~h.

In \textit{sog1-1} the same genes are flat, with only
\NumSentinelsRespondingMutant{} reaching 2-fold and RAD51 falling from $5.95$ to $0.60$.
The pipeline therefore recovers the central genetic finding of
\citet{bourbousse2018sog1} --- that the response is SOG1-dependent --- before any
regulatory model is fitted. This is the strongest available evidence that the
cross-study assembly is measuring biology rather than batch structure, because no step in
the harmonisation has access to genotype.

\begin{figure}[t]
\centering
\includegraphics[width=\linewidth]{../figures/fig2_sentinel_kinetics}
\caption{DNA-damage sentinel kinetics after cross-study harmonisation. All
\NumSentinelsResponding{} sentinels exceed 2-fold induction in wild type and trace a
rise--peak--decay curve; the same genes are flat in \textit{sog1-1}. No step in the
harmonisation has access to genotype.}
\label{fig:sentinels}
\end{figure}

The held-out OSD-782 arm shows no sentinel induction (maximum $|\log_2| = 1.0$). This is
the expected biology, not a failure: at 10--100~cGy the dose is two to three orders of
magnitude below that of the core experiment, and the source study reports ethylene
signalling rather than a canonical DNA-damage response. The sentinel gate is therefore
enforced only on the wild-type primary arm.

\subsection{The prior is incomplete in exactly the place that matters}

The DAP-seq prior covers \NumPriorTFs{} TFs and \NumPriorEdges{} edges, but does not
contain SOG1: DAP-seq never assayed it. A DREM model built on DAP-seq alone therefore
cannot attribute any bifurcation to the master regulator of the response being modelled,
and would do so silently. Calling SOG1 peaks from the same experiment's ChIP-seq recovers
\NumSogOneTargets{} target genes and closes the gap. We report this because the failure
mode generalises: a regulatory prior's absences are invisible in its outputs, and any
DREM analysis should check that the regulators its conclusions depend on are actually
present in the prior.

\subsection{Cell-type weighting}

Deconvolution places the irradiation cohort predominantly in rosette and seedling
clusters, consistent with whole-seedling material.
\NumWeightedEdgesInformed{} edges (\PctWeightedEdgesInformed{}\%) are atlas-informed and
\PctWeightedEdgesChanged{}\% change by more than $0.01$ under weighting. The
auto-decoder's \NumLatentDims{}-dimensional latent code moves monotonically away from the
untreated state over pseudo-time (Figure~\ref{fig:latent}).

\begin{figure}[t]
\centering
\includegraphics[width=0.72\linewidth]{../figures/fig3_latent_trajectory}
\caption{Displacement of the auto-decoder latent code from the earliest timepoint, per
arm. One interpretable scalar summarising the \NumLatentDims{}-dimensional trajectory.}
\label{fig:latent}
\end{figure}

Across \NumAblationSplits{} bifurcations in the wild-type model, compared on the
\AblationTestType{} test all three priors produced,
\NumAblationSplitsChanged{} (\PctAblationSplitsChanged{}\%) change their top-ranked TF
between the flat, binding and cell-type-weighted priors.

The split is not random. The bifurcations whose attribution is \emph{stable} across all
three priors are the biologically canonical ones: two are attributed to SOG1
(AT1G25580) and one to MYB3R1 (AT4G32730) --- the activator and repressor arms the source
study identified. The bifurcations that \emph{move} are the weaker ones, reassigned among
MYB-related, WRKY, TCP, NAC and bZIP regulators (WRKY33, TCP15, ANAC071, bZIP3, MYB49).
The prior therefore does not overturn well-supported attributions; it changes the calls
that were marginal to begin with, which is the behaviour a prior ought to have.

The wild-type arm is the only one with enough resolved bifurcations to read an ablation
from. The remaining three arms each resolved a single split shared by all three priors,
so their apparent 100\% change rate rests on one comparison and is reported in
Figure~\ref{fig:ablation} with its denominator rather than as a rate.

\begin{figure}[t]
\centering
\includegraphics[width=0.8\linewidth]{../figures/fig4_prior_ablation}
\caption{Proportion of bifurcations whose top-ranked transcription factor changes with
the prior. Labels give changed/total splits; pale bars rest on fewer than four
bifurcations and should not be read as rates.}
\label{fig:ablation}
\end{figure}

\subsection{The genotype contrast recovers the published regulatory logic}

The cohort's built-in positive control is decisive. Under all three priors, SOG1 ranks
\textbf{first} of the \NumPriorTFs{} scored TFs in the wild-type model. In
\textit{sog1-1} it collapses --- to rank \SogOneRankMutFlat{} under the flat prior and
\SogOneRankMutWeighted{} under the cell-type-weighted prior (Table~\ref{tab:genotype}).
MYB3R1 ranks \MybThreeRRankWTFlat{} in wild type and is retained in the mutant, as
expected for a repressor whose activity does not depend on SOG1.

This reproduces the regulatory logic \citet{bourbousse2018sog1} reported --- SOG1 as the
major activator, MYB3R as the major repressors --- from an independently reconstructed
pipeline, and it does so through the TF attribution rather than through expression alone.
Note that the attribution is only possible because SOG1 edges were added from ChIP-seq:
with the DAP-seq prior alone, SOG1 has no edges and could not have been ranked at all.

\begin{table}[t]
\centering
\caption{SOG1 rank among \NumPriorTFs{} scored transcription factors. Lower is more
significant; the mutant is the negative control.}
\label{tab:genotype}
\begin{tabular}{lrr}
\toprule
Prior & Wild type & \textit{sog1-1} \\
\midrule
flat     & \SogOneRankWTFlat{}     & \SogOneRankMutFlat{} \\
weighted & \SogOneRankWTWeighted{} & \SogOneRankMutWeighted{} \\
\bottomrule
\end{tabular}
\end{table}

\subsection{Independent corroboration from a different method}

The sibling analysis in \texttt{Plant\_response\_to\_radiation} models an overlapping set
of OSDR studies with a Gaussian-process autoencoder and WGCNA co-expression modules, with
no DREM and no TF attribution anywhere in it. Over the
\NumSharedGenesWithSibling{} genes both analyses modelled,
\NumNodesEnrichingModule{} of \NumSiblingNodes{} DREM nodes enrich one of the
\NumSiblingModules{} WGCNA modules at Bonferroni-corrected $p < 0.05$. The two methods
share raw data but almost no methodology, so the agreement is external corroboration
rather than a restatement.

\subsection{What the cell-type weighting actually did}

The weighting is a multiplier in $[0,1]$, so every regulator's total edge mass shrinks;
what re-ranks regulators against one another is shrinking more or less than the typical
TF. Measured that way, the atlas informs \NumTFsAtlasInformed{} of \NumTFsTotal{} TFs and
its action is specific rather than diffuse (Figure~\ref{fig:tissue}).

The six most strongly demoted regulators are \textit{FUS3}, ANAC079, ANAC058, ANAC038 and
ANAC029, together with AT1G04880 --- and every one of them has its expression maximum in a
silique or developing-seed cluster. The Bourbousse experiment irradiated \emph{seedlings}.
DAP-seq reports that these factors can bind their motifs, but they are not present in the
tissue being assayed, and the atlas removes them from contention. The regulators promoted
in relative terms --- TCP1, LEP, STZ, FAR1, WRKY59, ERF4 --- peak instead in rosette and
15-day seedling clusters, which is the tissue actually in the tube.

This is worth stating precisely because SOG1 is itself a NAC (ANAC008). The lever is not
down-weighting a family; it is discriminating \emph{within} one, demoting the seed-restricted
NACs while leaving the seedling-expressed ones untouched. That is also why the ablation
left the SOG1 and MYB3R1 attributions unmoved while reassigning the marginal calls: the
prior only had licence to act where the atlas disagreed with the binding data.

\subsection{Tissue association of the response}

Where the response is expressed was tested directly against the atlas rather than through
the deconvolution, for a reason worth recording. The per-timepoint NNLS fractions are not
stable: the median relative swing between adjacent timepoints is \ClusterSwing{} at cluster
resolution, and aggregating to organ level does \emph{not} fix it (\OrganSwing{}), because
the fit flips between clusters of the same organ as readily as across organs. The
time-averaged composition that the weighting consumes is stable (bootstrap CV
\TimeAveragedCV{}), but no per-timepoint fraction in this study should be read as
compositional dynamics, and none is.

Testing the gene sets against the atlas expression matrix directly --- no deconvolution ---
gives a clear split. The MYB3R-repressed arm is strongly and coherently tissue-restricted,
reaching $z = 26.2$ in silique, $23.3$ in 6-day seedling and $23.2$ in 21-day rosette
clusters, with 35 of 183 clusters significant. The SOG1-dependent arm is far more diffuse:
20 of 183 clusters, topping out at $z = 8.7$.

The asymmetry is interpretable and, we think, expected. Only a dividing cell has a G2/M
programme to shut down, so the repressed arm can only be read where proliferation is
happening; any cell can mount a damage response, so the activated arm is broadly
expressed. The radiation response is therefore not uniformly tissue-restricted --- its
\emph{repressive} half is, and its \emph{activating} half is not.

\begin{figure}[t]
\centering
\includegraphics[width=\linewidth]{../figures/fig6_tissue_association}
\caption{Tissue attribution. \textbf{Left:} deconvolved organ composition per timepoint ---
shown to make the instability visible, not to be read as dynamics. \textbf{Centre:}
enrichment of the SOG1-dependent set across atlas clusters, computed from the atlas
expression matrix directly. \textbf{Right:} regulators promoted or demoted by the
weighting relative to the median TF, labelled with the cluster where each peaks; the
demoted set is dominated by silique-restricted NACs absent from irradiated seedlings.}
\label{fig:tissue}
\end{figure}

\subsection{A decoder built on the signature, applied across OSDR}

The wild-type/\textit{sog1-1} contrast defines gene sets no single-genotype experiment
could: \NumSignatureSogDep{} genes induced in wild type, bound by SOG1, and losing that
induction in the mutant, against \NumSignatureSogIndep{} induced SOG1 targets that respond
regardless of genotype, and \NumSignatureMybRep{} MYB3R-repressed genes. Scoring these
across every \textit{Arabidopsis} study in OSDR with an expression matrix gave
\NumDecoderContrasts{} contrasts from \NumDecoderStudies{} studies.

Scoring the activated arm alone is not sufficient: several DNA-damage-related mutants
raise it without any exposure. What identifies acute irradiation is the coordinated
pattern the source study described --- SOG1 targets up \emph{and} MYB3R targets down at the
same time --- so the index is the conjunction
$\min(z_{\text{SOG1-dep}},\, -z_{\text{MYB3R-rep}})$, requiring both. The difference of the
two arms was tried first and rejected: it promoted six spaceflight and altered-gravity
contrasts on the strength of the repressed arm alone, with SOG1-dependent $z$ near zero.

At a fixed threshold of $1.96$ --- each arm individually significant, not fitted to these
data --- the decoder recovers \NumPositivesDetected{} of \NumLabelledPositives{} labelled
irradiation contrasts with \NumFalsePositives{} false positives among
\NumOtherExposures{} other exposures (AUC \DecoderAUC{}), and its scores order by dose:
Co-60 $\gamma$ at $16.8$--$20.3$, mixed radiation $7.8$, 100~cGy $5.2$, and 10~cGy not
detected. The single miss is that 10~cGy arm, two orders of magnitude below the dose the
signature was derived at. The \textit{sog1-1} and \textit{myb3r135} genotype contrasts
score at the opposite extreme ($-20.3$, $-21.3$, $-13.2$), which is the decoder
recognising the knockouts as the inverse of the response it was built on.

\begin{figure}[t]
\centering
\includegraphics[width=\linewidth]{../figures/fig5_decoder}
\caption{The radiation decoder. \textbf{Left:} every contrast in the two-arm plane; the
shaded corner is the decision region, entered only when both arms clear threshold.
\textbf{Right:} the fourteen highest-scoring contrasts. Labelled irradiation in red.}
\label{fig:decoder}
\end{figure}

\subsection{No other OSDR plant study carries the signature}

\NumFalsePositives{} non-radiation contrasts reach the threshold. Spaceflight, altered
gravity, and every mutant and treatment contrast in the repository fail it. On this
evidence, and within the limits set out below, spaceflight in these datasets does not
elicit the acute SOG1-dependent DNA-damage transcriptional programme.

That is not the same as finding nothing. \NumProliferationOnly{} contrasts show a distinct
and consistent pattern: the MYB3R/G2/M arm repressed with the SOG1 arm flat --- across
spaceflight and altered-gravity contrasts, a mean MYB3R-arm $z$ of $5.19$ against a mean
SOG1-arm $z$ of $-0.07$. Read as proliferation slowing without DNA damage, it is
biologically unremarkable; read as a methodological result it is the whole argument for
the conjunctive index, because a decoder scoring the repressed arm alone would have
reported all of it as radiation.

\subsection{Comparison with the published model}

\citet{bourbousse2018sog1} reported \NumPublishedGroups{} co-expressed gene groups from
six timepoints of OSD-510. Our flat-prior wild-type model over the pooled, denser series
yields \NumWTNodes{} nodes across \NumWTSplits{} splits from
\NumWTGenesAssigned{} assigned genes. These counts are not expected to be identical: the
pooled series adds the 10 and 45~min timepoints and draws on three accessions rather than
one, and DREM's node count responds to both. The comparison is reported for calibration.

\begin{table}[t]
\centering
\caption{Cohort composition, measured from the OSDR metadata API.}
\label{tab:cohorts}
\begin{tabular}{llrl}
\toprule
Cohort & Studies & Timepoints & Span \\
\midrule
A ($\gamma$-irradiation) & \NumCohortAStudies{} & \NumCohortATimepoints{} & 10 min -- 72 h \\
B (spaceflight)          & \NumCohortBStudies{} & \NumCohortBTimepoints{} & 4 -- 8 d \\
\bottomrule
\end{tabular}
\end{table}
